\hypertarget{appendix-b081b7abe193}{}
\section{Call stack and stack frames}\label{appendix-b081b7abe193-call-stack-and-stack-frames}

\textbf{Curriculum depth:} Introduced · \textbf{First seen:} C003 · \textbf{Latest development:} C003

\subsection{Definition and mechanism}\label{appendix-b081b7abe193-definition-and-mechanism}

A call stack conventionally holds LIFO activation records: return state, saved registers, spills, some arguments/locals, and alignment. Optimized locals may use registers or disappear; automatic storage does not mandate a stack slot. Downward growth is platform convention. Capacity/guard mappings can limit recursion despite free RAM elsewhere.

\subsection{Prerequisites}\label{appendix-b081b7abe193-prerequisites}

\begin{itemize}
\tightlist
\item
  \href{../05-Engine-I/main.pdf\#nameddest=appendix-0b104bcba06b}{Storage duration and object lifetime (05-Engine-I, appendix-0b104bcba06b)} --- Access requires an object and backing storage that are still valid.
\item
  \hyperref[appendix-c2a9ad1e2b49]{Process isolation and page permissions} --- The mapping context and permissions constrain which accesses proceed.
\end{itemize}

\subsection{Used by}\label{appendix-b081b7abe193-used-by}

Update downstream uses when later chunks build on this concept.

\subsection{Recall and next depth}\label{appendix-b081b7abe193-recall-and-next-depth}

Trace A→B→C→returns and identify the dead local.

\textbf{Pending:} Calling conventions, unwinding, coroutines, guard regions.

\subsection{Source and navigation}\label{appendix-b081b7abe193-source-and-navigation}

\hyperref[01-foundations-i-ch002-where-the-bytes-live]{C003 teaching chapter} · \hyperref[appendix-fd090d982a15-m003---parts-41-60]{Part definitions} · \href{../02-Foundations-II/main.pdf\#nameddest=appendix-4ad47bdd269a}{Dependency Map (02-Foundations-II, appendix-4ad47bdd269a)}
